\section{Encoding boundaries and finite checks}
\label{sec:checks}

\subsection{Precision and representation boundaries}

At $S=\Theta(n)$, the high-precision error in
Theorem~\ref{thm:capp} is still much larger than $2^{-n}$. Its
absolute-error guarantee does not distinguish zero accepting inputs
from one. Exact modes require their weight-cost budget or a supplied
short factorization. A general polynomial-bit equality constraint on
$n$ input bits need not satisfy those exact budgets.

The sum in~\eqref{eq:budget} is a sufficient condition for the
constructions above; no necessity is claimed for circuits whose
gates share realizations. The seed estimate used by
the algorithm is $R=O_{a,c}(S^{9/10})$, measured in the scale $S$.
Exact-only assignments have $R=0$; hashed assignments use their
pointwise error bounds for the precision comparison.

A ring calculation clarifies the cost of a particular congruence
filter. In $R=(\Z/2^P\Z)[V]/(V^\ell)$ let $\mathfrak m=(2,V)$.
If units $\zeta_j\in R$ and coefficients $c_j\in R$ satisfy
\[
 \sum_jc_j\zeta_j^u=2^e[u=0],\qquad 0\le u<M,\quad e<P,
\]
then $e\ge M-1$. Every unit is one modulo $\mathfrak m$; the
$(M-1)$st forward difference of the left side belongs to
$\mathfrak m^{M-1}$, while that of the right side is
$(-1)^{M-1}2^e$. The constant coefficient of
$\mathfrak m^{M-1}$ is divisible by $2^{M-1}$, forcing the bound.
Filters of this specified unit-power form therefore consume precision
proportional to their modulus range.

\subsection{The fixed-width counting interface}

The local condition ``two selected incident edges at every vertex''
recognizes a spanning 2-factor. Two disjoint triangles satisfy it and
have no HC. Joining them by a bridge and leaving the bridge unselected
gives the same example on a connected input. Rooted consistent cuts
give exactly $2^{\cc(F)-1}$ extensions and filter disconnected
2-factors modulo two. Distinct Hamiltonian solutions in one weight
layer can still cancel. The deterministic proofs use nonnegative
determinant sums instead.

Section~\ref{sec:lean-connectivity} supplies a linear fixed-width encoding
of each Cut\&Count weight layer, giving randomized decision and witnesses
under \HLean. Numbering-based certificates with $\Theta(n\log n)$
Boolean variables would give only $2^{O(n\sqrt{\log n})}$ even under
the manuscript's linear-density bound. The digit grouping avoids that
cost for fixed polynomially bounded coefficients. The deterministic
decision, full integer counting, and long-cost results of
Sections~\ref{sec:determinants}--\ref{sec:directed-routes} still use \HTC;
the new reduction does not establish those statements from \HLean.

One \HBB-only example can be verified directly. On a sparse graph,
exact chromatic number and an optimal coloring are computable,
conditional on \HBB, in
$\poly(n+L)2^{O(n\log\log n/\sqrt{\log n})}$ time and polynomial
space. Take $D=\lceil\log_2(n+2)\rceil^2$ and let $H$ contain
vertices of original degree greater than $D$. Then
$|H|=O(n/\log^2n)$. Compute $\chi(G[H])$ by enumerating all
$|H|^{|H|}$ color labelings. If it is at least $D+1$, an optimal
coloring of $H$ extends greedily to the remaining vertices and is
optimal. Otherwise try $q\le D+1$ and every $q^{|H|}$ coloring of
$H$. Represent each remaining color in $O(\log q)$ bits, reject
codes at least $q$, and test each remaining or crossing edge with an
$O(\log q)$-gate inequality circuit. The functional Tseitin 3-CNF
has $O(n\log\log n)$ variables and clauses. Add one fixed-zero dummy
bit per clause, plus twenty, to make density less than two while
preserving the count. Apply \HBB\ with fixed $k=3$ and density two;
all enumeration factors are absorbed in the stated exponent. Fixing
original color bits recovers an optimal extension. Empty graphs and
the case $q=1$ are checked directly. This is a routine
composition of degree splitting, binary comparison, and \HBB.

\subsection{Budget scales in lower-bound simulations}

All \HTC\ calls in Section~\ref{sec:lower} go through
Theorem~\ref{thm:capp}, whose scale is the actual $S$ of the query
circuit; the proofs use only upper bounds on the resulting running time,
so no lower bound on the size of a guessed circuit is needed.

For average-case amplification, the input length grows by $k$ and the
analysis-query exponent is proportional to $ck$. Condition
\eqref{eq:avg-time-condition} permits constant $k$ at a fixed positive
wire constant. Increasing CAPP precision alone does not remove that
time constraint.

\subsection{Classical sparse baselines}

The sparse reductions for the usual ETH comparison have explicit
linear budgets. Sparsification~\cite{IPZ01} gives a 3-CNF with $v$
variables and $a=O(v)$ clauses. The standard vertex-cover construction
has a variable edge and a clause triangle, with each clause vertex
joined to its literal's variable vertex. Its budget is $v+2a$;
there are $2v+3a$ vertices and at most $v+6a$ edges. The forced
one-per-variable and two-per-triangle choices make an unchosen clause
vertex imply a true literal. The standard dominating-set construction
has triangles $(t_i,f_i,z_i)$ and a vertex for each clause joined to
its literals, with budget $v$. The private $z_i$ forces one selection
per triangle; choosing $z_i$ can be replaced by a literal. It has
$3v+a$ vertices and at most $3v+3a$ edges.

From vertex cover, make a new triangle on each retained edge by adding
one vertex. Deleting a cover leaves a forest; a selected new vertex
in an FVS can be replaced by an endpoint without creating a cycle.
Adding a universal root and $k+2$ private leaves gives connected
vertex cover with budget $k+1$. For a dominating-set instance use
selectors $a_u$, elements $b_v$, edges $a_ub_v$ for $v\in N[u]$,
and a root joined to every selector. Root plus $k+2$ private leaves,
with budget $k+1$, gives connected domination. Selected elements can
be replaced by adjacent selectors because the root already dominates
all selectors. Alternatively declare the root and all elements Steiner
terminals: a tree using $s$ selectors has $n+s$ edges, so budget
$n+k$ is equivalent to $s\le k$. These graphs have linear vertex
and edge counts, and edge-to-vertex ratios at most four. They give
the standard sparse ETH baselines for these problems. The algorithms
in this paper use their explicit \HTC\ encodings.

\subsection{Ancillary computations}

The directory \ancfile{anc/} contains finite checks of the reductions.
Each check prints one summary line and writes its details to
\ancfile{anc/results/}. The counting checks use
integer or rational arithmetic; the floating-point threshold scan is
a diagnostic. Full exhaustive families and exhaustive
checks on sampled supports are distinguished in Table~\ref{tab:checks}.
The runner executes the checks serially.

\begin{table}[ht]
\centering\small
\renewcommand{\arraystretch}{1.12}
\begin{tabular}{@{}P{.37\linewidth}P{.55\linewidth}@{}}
\toprule
Check & Actual finite coverage\\
\midrule
\ancfile{finite_audit.cpp} & All permutations and interval partitions through $N=7$: 21,413,481 sign identities; 254,254 undirected marker cases through $n=5$; 333,544 directed cases through $n=4$\\
\ancfile{p6_sign_check.py} & 16,362 supported permutation identities; 350 sampled-support factor sums, including 218 explicit parity-wire sums\\
\ancfile{p7_p8_filter_check.py} & 120 sampled undirected graphs, 76 directed graphs; object sums and 68 metric-TSP comparisons\\
\ancfile{e2e_hc_network.py} & Nine graph/blocking combinations, including disjoint triangles; modular readback and integer division\\
\ancfile{check_gadgets.py} & 9,000 arithmetic replays of sliced threshold/equality/modular formulas, residue hashes, and small modular-function enumeration; no tensor contraction\\
\ancfile{check_G1_tensor.py} & 3,816 consistent-input factor contractions and 411 independently assigned slice-copy cases, including full auxiliary code domains\\
\ancfile{check_G3_tensor.py} & 136,896 residue-chain factor contractions on sampled instances, all their inputs and 24 primes; invalid residue codes included\\
\ancfile{protocol_exact_check.cpp} & Two parties, eight-bit inputs, 2,048-prime family: 3,346,110 nontrivial state/threshold classes\\
Explicit protocol tensor checks & Local tensors, equality subdivision, generic exact elimination, early stops and wrong-hash seeds; one extension for every original input\\
\bottomrule
\end{tabular}
\caption{Validation of the reductions. Auxiliary engine, block-count,
and unused-input checks and commands appear in \ancfile{anc/README.md}.}
\label{tab:checks}
\end{table}

The program \ancfile{hc_accounting.py} gives approximately $42.8S_G$ to
$52.4S_G$ bipartite wires on its five listed sparse fixtures, with
$P_0=4S_G+1$. These are finite examples; the asymptotic bounds follow
from the port counts proved above. The computations do not execute
the algorithm asserted by \HTC. Source hashes, commands, logs,
and scopes accompany the programs.
The arithmetic replay in \ancfile{check_gadgets.py} does not test
tensor-level extension counts, invalid auxiliary codes, or independently
assigned copies; its modular enumeration does not implement the general
integer-representation construction.
The tensor checks \ancfile{check_G1_tensor.py} (sliced gadgets) and
\ancfile{check_G3_tensor.py} (residue chains) cover the stated full-code
domains. In the sliced-gadget check, independently assigned slice copies test the local gadget's
extension count; the circuit's equality trees impose consistency of
those copies. The ancillary numerical threshold scan is a finite scan,
and the eventual inequalities are justified by the asymptotic proofs.

For Section~\ref{sec:lean-connectivity}, \ancfile{cnf_circuits.py}
constructs the functional 3-CNF gates and grouped weighted sums.
The \ancfile{check_cnf_*.py} programs compare literal-level exact
counts with direct graph and arithmetic references. They test
weight-layer cut multiplicities, odd-branch recovery, polynomial-cost
optimization, marker isolation, terminal-selection units, and the
unsigned modular expansion. Counterexamples include a feasible but
suboptimal TSP output and a missing-terminal directed Steiner formula.
The arithmetic checks include fixed inputs across a radix boundary;
large gate-accounting cases are distinguished from actual CNF solving.
The precise graph ranges, commands, and elapsed times are recorded in
\ancfile{anc/README.md} and \ancfile{anc/results/manifest.json}.
These checks use elementary exact counters, not the fast algorithm
assumed by \HLean, and do not formalize the reductions in Lean.

\noindent\begin{minipage}{\linewidth}
\subsection*{Use of AI}
The proofs and the text were produced with AI systems (GPT-6.1 Sol,
GPT-6 Astra, Claude Opus 5.5, and Claude Fable 5.1) under the author's
direction. They were checked by independent reviews with AI models and by the
computations in the ancillary files. No human expert has reviewed
them. The author takes responsibility for the content.
\end{minipage}\par
